\subsection{The regularization path}
\label{sec:regularization-path}

This subsection proves Theorem~\ref{thm:regularization-path} and then refines
it with the mass and spatial-mark limits in
Proposition~\ref{prop:path-masses}. The fixed-parameter support identities can
be made uniform across a compact part of the critical window. To obtain a path
limit, we must also resolve what happens at each activation threshold: as the
variance defect passes the threshold associated with one radial observation,
the corresponding edge mass reaches zero transversally. We first record the
radial process that supplies these thresholds, then prove the required uniform
continuation.

For dimension two and dimensions at least three, respectively, set
\begin{equation}
\label{eq:path-radial-processes}
 \mathcal X_n^{(2)}
 =\sum_{i=1}^n\delta_{\zeta_{n,i}^{(2)}},
 \quad
 \zeta_{n,i}^{(2)}=\frac{|Z_i|^2}{2}-\log n,
 \qquad
 \mathcal X_n^{(\ge3)}
 =\sum_{i=1}^n\delta_{\zeta_{n,i}^{(\ge3)}},
 \quad
 \zeta_{n,i}^{(\ge3)}
 =\frac{|Z_i-\bar Z_n|^2/(2\sigma_n^2)-b_n}{a_n}.
\end{equation}
Thus the exact excess in \eqref{eq:exact-transition-coordinate}, evaluated
along \eqref{eq:intro-highd-path-parametrization}, is
\(\xi_{n,i}^\star(s)=\zeta_{n,i}^{(\ge3)}-s\).

The one-level exceedance limits used to prove
Theorem~\ref{thm:main-support-law} hold jointly at every finite collection of
levels. Equivalently, the whole radial cloud has the following limit.

\begin{proposition}
\label{prop:path-radial-process}
In dimension two, and for every sequence \(d_n\ge3\),
\begin{equation}
\label{eq:path-radial-ppp}
 \mathcal X_n^{(2)}\Longrightarrow\mathcal X,
 \qquad
 \mathcal X_n^{(\ge3)}\Longrightarrow\mathcal X
\end{equation}
vaguely on \(\R\), where \(\mathcal X\) has intensity \(e^{-x}\,dx\).
For each fixed \(d\ge2\), the point-process convergence remains valid when
each radial excess is marked by the direction of its observation; the
limiting directions are independent and uniform on \(S^{d-1}\).
\end{proposition}

\begin{proof}
For \(d=2\), the variables \(|Z_i|^2/2\) are standard exponential, and hence
\(n\Pbb\{\zeta_{n,i}^{(2)}>x\}=e^{-x}\) for every fixed \(x\). Counts in
disjoint bounded intervals are multinomial rare-event counts, which proves the
first convergence.

For uncentered Gamma energies, Lemma~\ref{lem:gamma-hazard} gives, uniformly
for bounded \(x\),
\[
 n\overline F_n(b_n+a_nx)\longrightarrow e^{-x}.
\]
The rare-event multinomial calculation from the \(d=2\) case proves
convergence of the uncentered point process. In the range where empirical
centering is negligible on the
\(a_n\)-scale, Lemmas~\ref{lem:centered-radial-max} and
\ref{lem:centered-radial-max-high-d} transfer the process. In the remaining
range, the proof of Proposition~\ref{prop:direct-centered-gamma-max} gives
joint upper-tail factorization uniformly when finitely many thresholds lie in
\(b_n+a_nO(1)\). Although \eqref{eq:centered-joint-tail-factorization} is
displayed there with a common threshold, its determinant-and-tilting proof
also allows bounded coordinate-dependent offsets: use the corresponding
diagonal matrix of tilt parameters in place of the scalar tilt. The
determinant comparison and tilted local limit are unchanged.
Inclusion--exclusion then gives the factorization
for disjoint intervals, so the factorial-moment criterion proves the second
convergence in \eqref{eq:path-radial-ppp}. For fixed dimension, radius and
direction are independent; the marking theorem completes the proof.
\end{proof}

We next follow the optimizer while one point crosses its activation level. In
dimension two, let \(a_\lambda>0\) be the unique solution of
\(\Phi(a_\lambda)+\phi(a_\lambda)/a_\lambda=e^\lambda\), and put
\begin{equation}
\label{eq:path-d2-mass-profile}
 b_\lambda=1-e^{-\lambda}\Phi(a_\lambda),
 \qquad
 q_{\lambda,x}
 =\frac{[1-e^{-(x-x_\lambda)}]_+}{b_\lambda},
\end{equation}
and recall that \(\Lambda\) is the inverse of
\(\lambda\mapsto x_\lambda\). For \(d_n\ge3\), define
\begin{equation}
\label{eq:path-highd-scales}
 \gamma_n(s)=\frac{\sigma_n^2(b_n+a_ns)}{R_n^\star},
 \qquad
 \delta_n(s)=\frac{\sigma_n^2a_n}{\gamma_n(s)}.
\end{equation}

The next proposition is the uniform edge reduction. Its separation
hypothesis is only a proof device; Proposition~\ref{prop:path-radial-process}
will remove it.

\begin{proposition}
\label{prop:path-uniform-activation}
Fix compact intervals
\(J\subset\operatorname{int}(J^+)\Subset(0,\infty)\) and
\(I\subset\operatorname{int}(I^+)\subset\R\), and constants
\(C,M<\infty\) and \(\eta>0\). In dimension two, work on the event
\begin{equation}
\label{eq:path-d2-buffer-event}
 \max_i\zeta_{n,i}^{(2)}\le C,
 \qquad
 \#\{i:\zeta_{n,i}^{(2)}\ge x_{\inf J^+}\}\le M,
\end{equation}
on which the values \(\Lambda(\zeta_{n,i}^{(2)})\in J^+\) are pairwise at
least \(4\eta\) apart and are at least \(4\eta\) from the endpoints of
\(J^+\). Apart from an event whose probability tends to zero, each such point
has one activation value
\(\wh\lambda_{n,i}\), and
\begin{equation}
\label{eq:path-d2-activation}
 \max_i|\wh\lambda_{n,i}-\Lambda(\zeta_{n,i}^{(2)})|=o_{\Pbb}(1).
\end{equation}
Its patch contains one fitted atom when
\(\lambda<\wh\lambda_{n,i}\) and none when
\(\lambda\ge\wh\lambda_{n,i}\), and, uniformly for \(\lambda\in J\),
\begin{equation}
\label{eq:path-d2-patch-mass}
 n\pi_{n,i}^{(2)}(\lambda)
 =q_{\lambda,\zeta_{n,i}^{(2)}}+o_{\Pbb}(1).
\end{equation}
Points above \(x_{\sup J^+}\) are active throughout \(J\), points below
\(x_{\inf J^+}\) are inactive throughout \(J\), and the mass estimate applies
to every point satisfying the second cutoff in
\eqref{eq:path-d2-buffer-event}.

For every sequence \(d_n\ge3\), work instead on the event
\begin{equation}
\label{eq:path-highd-buffer-event}
 \max_i\zeta_{n,i}^{(\ge3)}\le C,
 \qquad
 \#\{i:\zeta_{n,i}^{(\ge3)}\ge\inf I^+-1\}\le M,
\end{equation}
on which the points in \(I^+\) are pairwise at least \(4\eta\) apart and at
least \(4\eta\) from the endpoints of \(I^+\). Apart from an event whose
probability tends to zero, each point in \(I^+\) has one activation value
\(\wh s_{n,i}\), and
\begin{equation}
\label{eq:path-highd-activation}
 \max_i|\wh s_{n,i}-\zeta_{n,i}^{(\ge3)}|=o_{\Pbb}(1).
\end{equation}
Its patch contains one fitted atom when \(s<\wh s_{n,i}\) and none when
\(s\ge\wh s_{n,i}\). Fixing \(s_0\in I\) and writing
\(\bar\delta_n=1\wedge\delta_n(s_0)\), its mass satisfies, uniformly for
\(s\in I\),
\begin{equation}
\label{eq:path-highd-patch-mass}
 n\pi_{n,i}^{(\ge3)}(s)
 =\left[1-e^{-\delta_n(s)(\zeta_{n,i}^{(\ge3)}-s)}\right]_+
  +o_{\Pbb}(\bar\delta_n),
\end{equation}
with the exponential expression replaced by its linearization when
\(d_nL_n/n^2\) is bounded below, the regime in which the centered cloud is
asymptotically simplex-like. Points above \(\sup I^+\) are active throughout \(I\), points
below \(\inf I^+\) are inactive throughout \(I\), and the mass estimate applies
to every point in the second cutoff in \eqref{eq:path-highd-buffer-event}. In
both dimensions there are no other noncentral contacts, and the NPMLE is
unique throughout the original compact interval.
\end{proposition}

\begin{proof}
We first explain why the fixed-parameter estimates are uniform. In dimension
two, the local diffuse background and the response of a radial point of excess
\(x\) are
\[
 B_\lambda(r)=e^{-\lambda}\Phi(-r),
 \qquad
 C_{\lambda,x}(r,z)=e^{-\lambda}e^{x-r^2/2-z^2/2}.
\]
Their spatial derivatives through order two and their first
\(\lambda\)-derivatives are bounded uniformly when \(\lambda\) ranges over a
compact subset of \((0,\infty)\). At finite \(n\),
\[
 \partial_\lambda e^{-\kappa_n|t|^2}
 =-\frac{|t|^2}{\rho_n^2}e^{-\kappa_n|t|^2},
 \qquad \kappa_n=\lambda/\rho_n^2,
\]
so the Gaussian envelopes used for the undifferentiated field also dominate
its \(\lambda\)-derivative on \(|t|=\rho_n+O(1)\). Applying the compact decomposition in
Proposition~\ref{prop:two-d-compact-decomposition}, the lower-cloud bounds, and
the \(C^2\) estimates of Section~\ref{sec:higher-dimensional-support-counts}
on a deterministic \(\lambda\)-grid, then using these derivative bounds
between grid points, makes every fixed-parameter estimate uniform in
\(C^2(r,z)\cap C^1(\lambda)\). The radial compactness and outside-window gaps
in Proposition~\ref{prop:d-two-radial-compactness} and
Lemmas~\ref{lem:d-two-diffuse-outside}--\ref{lem:d-two-tail-outside} are
uniform for the same reason: on \(J^+\), every damping constant stays in a
fixed compact subset of \((0,\infty)\).

It remains to inspect the limiting patch program
\eqref{eq:support-d2-local-program}. Its unique spatial optimizer is
\((-a_\lambda,0)\). After the other patch variables have been optimized, the
derivative at zero rescaled mass is
\begin{equation}
\label{eq:path-d2-boundary-score}
 S_{\lambda,x}(0)=b_\lambda\{e^{x-x_\lambda}-1\}.
\end{equation}
The functions \(a_\lambda,x_\lambda,b_\lambda\) are smooth on compact
intervals, with \(x_\lambda'>0\) and \(b_\lambda>0\). At the unique zero
\(\lambda=\Lambda(x)\),
\[
 \partial_\lambda S_{\lambda,x}(0)=-b_\lambda x_\lambda',
\]
which is bounded away from zero. The spatial Hessian there is
\(-b_\lambda I_2\), and the mass curvature is strictly negative. The uniform
finite-patch expansion therefore gives one empirical zero satisfying
\eqref{eq:path-d2-activation}. Strict concavity places positive mass below
that zero and zero mass at and above it. Solving the limiting interior score
equation gives \eqref{eq:path-d2-patch-mass}. Continuity of the optimizer at
zero makes the estimate uniform through the crossing.

For \(d_n\ge3\), write
\[
 h_n(s)=\frac{\gamma_n'(s)}{\gamma_n(s)}
       =\frac{a_n}{b_n+a_ns}.
\]
If \(q_j(s)=y_j/\sqrt{\gamma_n(s)}\), differentiation of an edge kernel gives
\begin{align}
 &\left|\partial_s
 \exp\left\{(q_i(s)+z)\cdot q_j(s)-\frac{|q_i(s)+z|^2}{2}\right\}\right|
 \notag\\
 &\hspace{35mm}\le
 C h_n(s)\{\log n+|q_j(s)-q_i(s)-z|^2\}
 \exp\left\{(q_i(s)+z)\cdot q_j(s)-\frac{|q_i(s)+z|^2}{2}\right\}.
\label{eq:path-kernel-derivative-main}
\end{align}
Moreover \(h_n(s)\log n\asymp\delta_n(s)\), uniformly on \(I\). Thus the
order-zero and order-two absolute envelopes already proved in the spatial,
entropy, and support-recovery arguments also control one derivative in \(s\).
As in dimension two, a slowly refining deterministic grid makes all buffered
support identities uniform on \(I\).

Near one crossing, optimize the central contact and every other patch, leaving
only the rescaled mass \(m_i=n\pi_i\ge0\) in the crossing patch free. The
algebraic expansion in the proof of
Lemma~\ref{lem:uniform-finite-edge-expansion} is obtained before the threshold
buffer is used: its envelope assumptions do not depend on the sign of the
crossing excess. Holding this one mass free therefore gives the expansion in
Lemma~\ref{lem:uniform-finite-edge-expansion} uniformly while
\(|\zeta_{n,i}^{(\ge3)}-s|\le2\eta\), with every other patch
still buffered. In particular, the boundary score and its derivative are
\begin{align}
 \partial_{m_i}\mathcal L_{n,i}(s,0)
 &=e^{\delta_n(s)(\zeta_{n,i}^{(\ge3)}-s)}-1
   +o_{\Pbb}(\bar\delta_n),
\label{eq:path-highd-boundary-score}\\
 \partial_s\partial_{m_i}\mathcal L_{n,i}(s,0)
 &=-\delta_n(s)+o_{\Pbb}(\bar\delta_n),
\label{eq:path-highd-boundary-derivative}
\end{align}
while its second mass derivative is \(-1+o_{\Pbb}(1)\). Indeed,
\(\delta_n'=-h_n\delta_n\), so differentiating the leading exponent adds
only \(O\{h_n\delta_n|\zeta_{n,i}^{(\ge3)}-s|\}\) to
\(-\delta_n\); inequality \eqref{eq:path-kernel-derivative-main} keeps every
remainder \(o_{\Pbb}(\bar\delta_n)\) after differentiation. When
\(d_nL_n/n^2\) is bounded below, the ordered-contact argument gives
\eqref{eq:path-highd-boundary-score}--\eqref{eq:path-highd-boundary-derivative}
with the exponential linearized. More explicitly, the ordered-contact
response remainder in
\eqref{eq:potts-corrected-height} is bounded by a constant times
\[
 \delta_n^2+\frac{C_n^2}{n}
 +\frac{\log n\,(C_n+\delta_n)^2}{n^2}.
\]
Here \(\delta_n'=-h_n\delta_n\), and the Gram-error array defining \(C_n\)
is rescaled by \(\gamma_n(s)^{-1}\), so its upper Dini derivative is at most
\(h_nC_n\). The rates in \eqref{eq:terminal-simplex-derived-rates} therefore
show directly that one \(s\)-derivative leaves this remainder
\(o_{\Pbb}(\delta_n)\).

Equation \eqref{eq:path-highd-boundary-derivative} is uniformly negative.
Hence each boundary score has exactly one zero, which proves
\eqref{eq:path-highd-activation}; strict mass curvature and the spatial
certificate give one contact below the zero and none at or above it. Solving
the interior score equation gives \eqref{eq:path-highd-patch-mass}. Finally,
the uniform residual KKT gaps proved in the dimension-specific support
identities exclude contacts outside these patches. The strict concavity of
the local mass programs and the uniqueness of each spatial contact then prove
simultaneous uniqueness, including at an activation value, where the crossing
patch has zero mass.
\end{proof}

\begin{proof}[Proof of Theorem~\ref{thm:regularization-path}]
Enlarge the parameter interval slightly. Proposition
\ref{prop:path-radial-process} makes the radial maximum and the number of
points in the enlarged critical window tight. The limiting process is simple,
so after imposing fixed radial and cardinality cutoffs, its activation levels
are separated from one another and from the enlarged interval's endpoints
with probability tending to one as the separation constant decreases to zero.
Apply Proposition~\ref{prop:path-uniform-activation} first with fixed cutoffs
and separation, then let the cutoffs tend to infinity and the separation tend
to zero. A diagonal choice gives a single event of probability tending to one
on which all of its conclusions hold.

In dimension two, the support-count path is the tail-count path of the
perturbed activation measure
\(\sum_i\delta_{\wh\lambda_{n,i}}\). Equation
\eqref{eq:path-d2-activation} allows a piecewise-linear time change that moves
each \(\wh\lambda_{n,i}\) to \(\Lambda(\zeta_{n,i}^{(2)})\) and converges
uniformly to the identity. Points above the enlarged interval contribute the
same constant to both paths. After adding them, the resulting path is exactly
\(\mathcal X_n^{(2)}((x_\lambda,\infty))\). Proposition
\ref{prop:path-radial-process} and the continuous-mapping theorem for tail
counts of a simple point process prove \eqref{eq:intro-d2-path-limit} in the
\(J_1\) topology.

In dimensions at least three, equation
\eqref{eq:path-highd-activation} supplies an analogous time change. After
that time change and the addition of
the points active throughout \(I\), the path is
\(\mathcal X_n^{(\ge3)}((s,\infty))\). This proves
\eqref{eq:intro-highd-path-limit}. Simultaneous uniqueness is part of
Proposition~\ref{prop:path-uniform-activation}.
\end{proof}

Proposition~\ref{prop:path-uniform-activation} also identifies the mass and
location carried by every jump. We state
the mass limit explicitly because it is stronger than the support-count path
and shows that the fitted mixing distribution itself changes continuously even
when its support size jumps.

\begin{proposition}
\label{prop:path-masses}
Associate each noncentral fitted atom with its radial observation, and extend
its mass by zero after the corresponding patch becomes inactive. In dimension
two, for compact \(J\Subset(0,\infty)\),
\begin{equation}
\label{eq:path-d2-mass-limit}
 \max_{i:\,\zeta_{n,i}^{(2)}\ge x_{\inf J}-1}
 \sup_{\lambda\in J}
 \left|n\pi_{n,i}^{(2)}(\lambda)
       -q_{\lambda,\zeta_{n,i}^{(2)}}\right|
 \xrightarrow{\Pbb}0,
\end{equation}
after first imposing and then removing a fixed upper radial cutoff. Moreover,
\begin{equation}
\label{eq:path-d2-total-mass-limit}
 \left\{n\sum_i\pi_{n,i}^{(2)}(\lambda):\lambda\in J\right\}
 \Longrightarrow
 \left\{\sum_{x\in\mathcal X}q_{\lambda,x}:\lambda\in J\right\}
 \qquad\text{in }C(J).
\end{equation}

For \(d_n\ge3\), fix \(s_0\in I\), pass to a subsequence on which
\(\delta_n(s_0)\to\delta\in[0,\infty]\), and define
\[
 \psi_\delta(y)=
 \begin{cases}
 (1-e^{-\delta y})/(1\wedge\delta),&0<\delta<\infty,\\
 y,&\delta=0,\\
 1,&\delta=\infty,
 \end{cases}
 \qquad y>0,
 \quad \psi_\delta(y)=0\quad(y\le0).
\]
Then
\begin{equation}
\label{eq:path-highd-mass-limit}
 \max_{i:\,\zeta_{n,i}^{(\ge3)}\ge\inf I-1}
 \sup_{s\in I}
 \left|\frac{n\pi_{n,i}^{(\ge3)}(s)}{\bar\delta_n}
       -\psi_\delta((\zeta_{n,i}^{(\ge3)}-s)_+)\right|
 \xrightarrow{\Pbb}0,
\end{equation}
and the sum of these profiles converges in \(C(I)\) to
\(s\mapsto\sum_{x\in\mathcal X}\psi_\delta((x-s)_+)\). For every fixed
\(d\ge2\), the convergence is joint with the radial directions; each active
atom has the direction of its associated extreme observation asymptotically.
In dimension two its radial score location is
\(\rho_n-a_\lambda+o_{\Pbb}(1)\), uniformly away from activation.
\end{proposition}

\begin{proof}
Equations \eqref{eq:path-d2-patch-mass} and
\eqref{eq:path-highd-patch-mass} give the individual limits on events with
fixed radial and cardinality cutoffs. Lemma~\ref{lem:gamma-hazard} gives
\(a_n/b_n\to0\), and hence
\(\delta_n(s)/\delta_n(s_0)\to1\) uniformly on \(I\); this turns
\eqref{eq:path-highd-patch-mass}, after division by \(\bar\delta_n\), into the
profile \(\psi_\delta\). Proposition
\ref{prop:path-radial-process} removes these cutoffs. The profiles are jointly
continuous in the path parameter and radial mark, including at activation
where they vanish. Only finitely many limiting Poisson points contribute on a
compact parameter interval. Summation is therefore a continuous map into
\(C(J)\) or \(C(I)\), which proves the total-mass limits. The marked part of
Proposition~\ref{prop:path-radial-process}, together with the unique spatial
optimizer in each local likelihood program, gives the final assertions.
\end{proof}
